\section{Related Work}
\label{sec:related-work}

\paragraph{Where defenses intervene.}
AgentDojo, ToolEmu, and AgentLAB expose indirect prompt injection in tool-rich
and long-horizon tasks~\cite{agentdojo,toolemu,agentlab}.  Step classifiers and
prompt defenses decide whether a proposed action appears safe
~\cite{toolsafe,safiron}; CaMeL and IPIGuard preserve trusted control or track
data dependencies~\cite{camel,ipiguard}.  AttriGuard and CausalArmor go further
by intervening on observations and rerunning the agent to estimate which input
caused a proposed call~\cite{attriguard,causalarmor}.  These approaches reason
upstream about why the planner selected an action.  Our registration procedure
intervenes downstream on arguments, defaults, and pre-state, executes the tool,
and asks which committed effects must remain distinguishable.  It produces an
object for authorization rather than another classifier of the planner's
choice.

\paragraph{What authorization sees.}
Agent policy systems bound tools, actions, parameters, or capabilities
~\cite{progent,miniscope,clawguard,scopegate,agentguard}.  PACT assigns semantic
roles to arguments, while AuthGraph compares clean authority with the
provenance of argument values~\cite{pact,authgraph}.  SecureClaw and commit-time
authorization emphasize checking fresh authority at the effect sink
~\cite{secureclaw,committimeauth}.  These systems are complementary to our
representation question.  An argument need not correspond one-to-one with an
effect: defaults and pre-state can create effects without an explicit value,
and one list argument can expand into several independently authorizable
changes.  We supply validated effect objects to a policy; we do not infer the
user's authority.  Contextual frameworks separately formalize task, action,
source, and data-isolation boundaries~\cite{siu2026formalizing}.  Our
enforcement boundary follows established complete
mediation, least privilege, confused-deputy, information-flow, and provenance
principles~\cite{saltzer1975protection,dennis1966programming,hardy1988confused,denning1976lattice,buneman2001why,cheney2009provenance}.

\paragraph{How effects are specified.}
Type-and-effect systems make program effects explicit by construction
~\cite{lucassen1988effects}; ETAS brings effect rows and compiled monitors to an
agent language~\cite{etas2026}.  Contract2Tool infers planning contracts,
ContractGuard protects supplied contracts, and ToolGuardian combines
descriptions, traces, mock execution, and source analysis to vet tools
~\cite{contract2tool,contractguard,toolguardian2026}.  We address a narrower
obligation that remains after an effect specification is proposed: does the
representation preserve every distinction required by the selected policy
family?  Paired executions can falsify that property for existing tools that
were not written in an effect-typed language.

\paragraph{How representations are validated.}
Counterexample-guided abstraction refinement starts with a coarse abstraction
and splits it when a concrete execution falsifies the abstraction
~\cite{clarke2000cegar}.  Causal abstraction similarly asks whether
interventions at a lower level are preserved by a higher-level model
~\cite{pearl2009causality,beckers2019abstracting,geiger2022inducing}.  We adapt
this idea to tool registration: a call or state intervention is useful only
when execution changes the committed effect, and it requires a representational
split only when the policy also separates the pair.  The conclusion is scoped
to the frozen intervention and policy domains, not unrestricted causal
discovery.
